% 第3章 集成运算放大电路
\chapter{集成运算放大电路}

\section{多级放大电路的一般问题}
\subsection{多级放大电路的耦合方式}
\begin{itemize}
    \item 直接耦合放大电路
    \item 阻容耦合放大电路
    \item 变压器耦合放大电路
    \item 光电耦合
\end{itemize}

\subsection{多级放大电路的动态分析}
多级放大电路的总电压放大倍数等于各级电压放大倍数的乘积，输入电阻为第一级的输入电阻，输出电阻为最后一级的输出电阻。

\section{集成运算放大电路概述}
\subsection{集成运放的电路结构特点}
集成运放采用直接耦合方式，具有高增益、高输入电阻、低输出电阻、宽频带等特点，适用于模拟信号的运算与处理。

\subsection{集成运放电路的组成及其各部分的作用}
\begin{itemize}
    \item 输入级：通常为差分放大电路，抑制零点漂移，提高输入电阻
    \item 中间级：电压放大级，提供足够大的电压放大倍数
    \item 输出级：互补对称输出电路，降低输出电阻，提高带负载能力
    \item 偏置电路：为各级电路提供稳定的静态工作电流
\end{itemize}

\subsection{集成运放的电压传输特性}
集成运放的电压传输特性分为线性区和非线性区：
\begin{itemize}
    \item 线性区：输出电压与输入差模电压成正比，$u_O = A_{od} \cdot u_{ID}$
    \item 非线性区：输出电压为正饱和值或负饱和值，$u_O = \pm U_{OM}$
\end{itemize}

\section{集成运放中的单元电路}
\subsection{直接耦合放大电路的零点漂移现象}
\subsubsection{零点漂移现象及其产生的原因}
零点漂移：输入信号为零时，输出信号随时间缓慢变化的现象。
主要原因：温度变化导致晶体管参数（$I_{CBO}$、$U_{BE}$、$\beta$）变化，电源电压波动等。

\subsection{差分放大电路}
\subsubsection{电路的组成}
由两个对称的单管放大电路组成，通过公共发射极电阻抑制共模信号，放大差模信号。

\subsubsection{长尾式差分放大电路}
\paragraph{静态分析}
静态时输入差模信号为零，由于电路对称，两管集电极电位相等，输出电压为零。

\paragraph{对共模信号的抑制作用}
共模信号是两管输入的相同信号，长尾电阻对共模信号产生强负反馈，抑制零点漂移。

\paragraph{对差模信号的放大作用}
差模电压放大倍数：
\[ A_{d} = -\frac{\beta \left( R_{c} \parallel \frac{R_{L}}{2} \right)}{R_{b} + r_{be}} \]
输入电阻：
\[ R_i=2(R_b+r_{be}) \]
输出电阻：
\[ R_o=2R_c \]

\paragraph{电压传输特性}
描述差分放大电路输出电压与输入差模电压之间的关系，分为线性区和饱和区。

\subsubsection{差分放大电路的四种接法}
根据输入、输出端的接地方式，分为：双端输入双端输出、双端输入单端输出、单端输入双端输出、单端输入单端输出。

\subsubsection{改进型差分放大电路}

\subsection{电流源电路}
\subsubsection{基本电流源电路}
\begin{itemize}
    \item 镜像电流源：$I_{C2} \approx I_R$，实现电流复制
    \item 比例电流源：
    \[ I_{C1}\approx \frac{R_{e0}}{R_{e1}} \cdot I_R \]
    \item 微电流源：产生微安级的恒定电流，适用于低功耗电路
\end{itemize}

\subsubsection{改进型电流源电路}
\begin{itemize}
    \item 加射极输出器的电流源：提高电流复制精度
    \item 威尔逊电流源：具有更高的输出电阻，电流稳定性更好
\end{itemize}

\subsubsection{多路电流源电路}
由一个基准电流源提供多路恒定电流，为集成运放各级电路供电。

\subsubsection{以电流源为有源负载的放大电路}
用电流源代替集电极电阻，提高电压放大倍数（电流源输出电阻大）。

\subsection{直接耦合互补输出级}
采用NPN和PNP管互补对称结构，消除交越失真，实现无变压器、无电容的直接耦合输出。

\section{集成运放电路简介}
\subsection{双极型集成运放电路}
以双极型晶体管为核心，具有高增益、高速度的特点，典型代表为$\mu$A741。

\section{集成运放的性能指标及低频等效电路}
\subsection{集成运放的主要性能指标}
\begin{itemize}
    \item 开环差模增益 $A_{od}$
    \item 输入失调电压 $U_{IO}$
    \item 输入失调电流 $I_{IO}$
    \item 输入偏置电流 $I_{IB}$
    \item 共模抑制比 $K_{CMR}$
    \item 带宽 $BW$
    \item 转换速率 $SR$
\end{itemize}

\subsection{集成运放的低频等效电路}
将集成运放等效为具有高输入电阻、低输出电阻、高开环增益的双端口网络，用于线性区的电路分析。

\section{例题参考}
2、5